\documentclass{beamer}

\usetheme{simple}

\usepackage{lmodern}
\usepackage[scale=2]{ccicons}

\usepackage[
    backend=biber,
    style=numeric,
    sorting=none
]{biblatex}

\addbibresource{references.bib}

% TODO: 
%   position adjustement
%   change colours
%       

% Watermark background (simple theme)

\setwatermark{\includegraphics[height=8cm]{img/Heckert_GNU_white.png}}


\title{Cpp26: Reflection}
\subtitle{}
\date{\today}
\author{Gautam Ghosh}
\institute{University of Stuttgart}

\begin{document}

\maketitle

\begin{frame}{}

  \begin{block}{}
  \begin{center}{\textbf{DISCLAIMER:}}\end{center}
  \begin{center}{AI tools were used to assist with the preparation of this presentation. It was used for the following tasks:
  \begin{itemize}
      \item Navigating documentation.
      \item Writing LaTeX code.
  \end{itemize}
    }\end{center}

  \begin{center}{All C++ code has been written and reviewed by the Author.}\end{center}
  \end{block}
  
\end{frame}

\begin{frame}{}

  \begin{block}{}
  \begin{center}{"There is no point in developing software unless you care about doing it well."}\end{center}

  \begin{center}{Andrew Hunt, David Thomas, The Pragmatic Programmer}\end{center}
  \end{block}
  
\end{frame}

\begin{frame}[plain]
    \centering
    \includegraphics[
        width=\textwidth,
        height=\paperheight,
        keepaspectratio
    ]{img/meme.png}
\end{frame}

\begin{frame}{The Reflection operator}

  \begin{itemize}
        \item Operator syntax: \textbf{\^\^}
        \item Proposed in the paper P2996R13.
        \item Encapsulates compile-time information about a C++ entity using the \textbf{std::meta::info} type.\textsuperscript{[2]}
        
    \end{itemize}
    \centering
        \includegraphics[
            width=\textwidth,
            height=0.6\textheight,
            keepaspectratio
        ]{img/reflection}
\end{frame}

\begin{frame}{The Splice operator}

    \begin{itemize}
        \item Operator syntax: \textbf{[:\textit{reflection}:]}
        \item Proposed in the paper P2996R13.
        \item Enables converting \textbf{std::meta::info} object into \textbf{source code}.\textsuperscript{[2]}
        \item Example:
        
    \end{itemize}
        \centering
        \includegraphics[
            width=\textwidth,
            height=0.85\textheight,
            keepaspectratio
        ]{img/splicer}
  
\end{frame}

\begin{frame}{The Expansion Statement}

    \begin{itemize}
        \item Proposed in the paper P1306R5.
        \item Enables unrolling for loops at compile time.
        \item The compound-statement does not need to be \textbf{constexpr}.
        \item Example:\textsuperscript{[1]}
        
    \end{itemize}
        \centering
        \includegraphics[
            width=\textwidth,
            height=0.85\textheight,
            keepaspectratio
        ]{img/expansion_statement}
\end{frame}

\begin{frame}{References}
    \footnotesize

    \begin{enumerate}
        \item \textbf{Expansion Statements}\newline
        \url{https://www.open-std.org/jtc1/sc22/wg21/docs/papers/2025/p1306r5.html#expansion-over-tuples}

        \item \textbf{Reflection for C++26}\newline
        \url{https://www.open-std.org/jtc1/sc22/wg21/docs/papers/2025/p2996r13.html#stdmetainfo}

        \item \textbf{cppreference}\newline
        \url{https://en.cppreference.com/cpp/header/meta}

        \item \textbf{C++26 reflection at compile-time}\newline
        \url{https://andreasfertig.com/blog/2025/08/cpp26-reflection-at-compile-time/}

    \end{enumerate}
\end{frame}

\begin{frame}[plain]
    \centering

    \vspace{1.5cm}

    {\Huge\textbf{Thank You!}}

    \vspace{0.8cm}

    {\Large Questions?}

    \vspace{1cm}

    {\normalsize
    Code available on GitHub:
    }

    \vspace{0.4cm}

    \includegraphics[
        width=4cm
    ]{img/cpp26_reflection_qr.png}

\end{frame}

\end{document}
